\section*{第 \ref{ch:g9:evidence-of-evolution} 章——演化的证据}

\begin{solution}{exo:g9:evidence-of-evolution:1}
一种保存在沉积岩中的过去生命的遗迹。
地层按沉积顺序逐层堆叠，最老的在最下面，
因此沿地层向上，就是沿时间向前。
\end{solution}

\begin{solution}{exo:g9:evidence-of-evolution:2}
\omterm{def:g1:animals-around-us:animal}{动物}群逐层变化，较深的地层中大多数\omterm{def:g5:biodiversity:species}{物种}已经\omterm{prop:g5:biodiversity:loss}{灭绝}；
各大类群按照系统树的嵌套次序依次出现——
\omterm{prop:g3:sorting-living-things:groups}{鱼类}在深层，\omterm{def:g1:my-body:parts}{四肢}类亲缘\omterm{def:g1:animals-around-us:animal}{动物}在其上，\omterm{prop:g3:sorting-living-things:groups}{哺乳动物}和\omterm{prop:g3:sorting-living-things:groups}{鸟类}位于最高处；
中间类型则位于\omterm{def:g1:plants-around-us:tree}{分支}点上。
\end{solution}

\begin{solution}{exo:g9:evidence-of-evolution:3}
\omterm{def:g1:animals-around-us:covering}{羽毛}和翅膀——即\omterm{prop:g3:sorting-living-things:groups}{鸟类}的组件——安装在
带有牙齿和长骨质尾巴的骨架上——即爬行类的组件：两套组件
同时出现在一种\omterm{def:g1:animals-around-us:animal}{动物}身上，正位于系统树所画的分叉处。
\end{solution}

\begin{solution}{exo:g9:evidence-of-evolution:4}
鲸体内隐藏的髋骨和腿骨：说明其祖先曾经行走。
不会飞的甲虫那被封闭的翅膀：说明其祖先曾经飞行。
人类的尾骨（或退化的阑尾）：说明祖先曾有尾巴，
且食性不同。每一项都是保留下来的部分，
见证了祖先的身体构造。
\end{solution}

\begin{solution}{exo:g9:evidence-of-evolution:5}
比较不同\omterm{def:g5:biodiversity:species}{物种}间同一\omterm{def:g9:genes-and-dna:gene}{基因}的碱基差异数：
差异越少，\omterm{prop:g6:classification-kinship:kinship}{亲缘关系}越近。将这些距离汇总，
即可构建一棵系统树——它与\omterm{def:g3:bones-and-muscles:skeleton}{骨骼}和
\omterm{def:g9:evidence-of-evolution:fossils}{化石}所绘制的系统树相吻合。
\end{solution}

\begin{solution}{exo:g9:evidence-of-evolution:6}
椒花蛾的深色型随烟灰与洁净空气的变化而变化；
抗杀虫剂和抗生素的害虫在数年内相继出现；
雀类喙的尺寸随干旱与
食物充裕而变化。
\end{solution}

\begin{solution}{exo:g9:evidence-of-evolution:7}
若重新设计，可以为每个肢体配备最佳机制——
如翅膀用支柱结构，鳍肢用水翼——
且不会重复那些任意的\omterm{def:g3:bones-and-muscles:skeleton}{骨骼}数目。唯一的排布，
生硬地被拉伸用于各种用途，正是“在\omterm{def:g9:heredity-and-traits:heredity}{遗传}基础上修补”所产生的结果，
而自由设计则不会产生这种结果。
\end{solution}

\begin{solution}{exo:g9:evidence-of-evolution:8}
类群特征：分泌\omterm{def:g2:animals-grow-up:born}{乳汁}、用肺\omterm{def:g7:what-is-respiration:respiration}{呼吸}、胎毛——\omterm{prop:g3:sorting-living-things:groups}{哺乳动物}的
\omterm{def:g9:heredity-and-traits:heredity}{遗传}印记（见\cref{ex:g6:classification-kinship:whale}）。
退化器官：隐藏的髋骨和腿骨。\omterm{def:g9:evidence-of-evolution:fossils}{化石}：
按时间排列的序列，从四足的滨岸\omterm{prop:g3:sorting-living-things:groups}{哺乳动物}
到无腿的游泳者。分子：其\omterm{def:g9:genes-and-dna:gene}{基因}距离把它置于
\omterm{prop:g3:sorting-living-things:groups}{哺乳动物}之中，与有蹄类亲缘最近。四方证据，一份
卷宗。
\end{solution}

\begin{solution}{exo:g9:evidence-of-evolution:9}
因为这些证据之间不可能串通：岩层的形成、
身体结构的形成以及分子序列的形成，
彼此互不相干，各自无法获知对方的证词。独立的证据
相互吻合，可成倍提高可信度——三类来源得出同一棵系统树，
正是该方法所看重的趋同性。
\end{solution}

\begin{solution}{exo:g9:evidence-of-evolution:10}
因为该论断禁止这种情况：\omterm{prop:g3:sorting-living-things:groups}{哺乳动物}不得出现在其祖先
所在地层之下。每一次挖掘都是新的证伪机会——
而历来所有发掘中持续存在的这种缺失，
正说明该论断在每一次
发掘中都在通过检验。
\end{solution}

\begin{solution}{exo:g9:evidence-of-evolution:11}
\omterm{def:g4:classifying-animals:vertebrate}{脊椎动物}\omterm{def:g5:human-reproduction:pregnancy}{胚胎}在早期
奠定体制的几周内高度相似，较晚时才出现分化。发育
一章中的机制解释了其原因：早期指令奠定了
共同的身体蓝图——即全书最深、最古老的“章节”——
而较晚的解读则写入该\omterm{def:g5:biodiversity:species}{物种}的
特有细节。早期文本相同，说明深层祖先相同。
\end{solution}

\begin{solution}{exo:g9:evidence-of-evolution:12}
论断：\omterm{def:g5:biodiversity:species}{物种}会发生变化，并拥有共同的祖先。预测：
\omterm{def:g1:animals-around-us:animal}{动物}群按时间排列并符合系统树的嵌套关系；中间类型
位于分叉处；存在祖先身体构造的遗迹；分子距离
与\omterm{prop:g6:classification-kinship:kinship}{亲缘关系}逐级对应。证据：岩石提供了
顺序与中间类型；身体结构提供了蓝图与
遗迹；分子提供了相吻合的系统树；活体
观察提供了正在发生的变化。结论：具有共同祖先的变化
得到独立证据趋同支持，
并为下一章提出了问题：\emph{驱动变化的
究竟是什么？}
\end{solution}

\begin{solution}{pb:g9:evidence-of-evolution:1}
\textbf{1.} 顺序：三叶虫位于最下层，其次是最早的\omterm{prop:g3:sorting-living-things:groups}{鱼类}、
两栖类亲缘\omterm{def:g1:animals-around-us:animal}{动物}的足迹、\omterm{prop:g4:classifying-animals:classes}{爬行动物}、\emph{Archaeopteryx}、
早期\omterm{prop:g3:sorting-living-things:groups}{哺乳动物}的颌骨。图注：“自下而上读岩层，
所呈现的变化，恰与系统树的预测
完全一致——中间类型都位于
\omterm{def:g1:plants-around-us:tree}{分支}点处。”

\textbf{2.} 图注：“同一套\omterm{def:g3:bones-and-muscles:skeleton}{骨骼}蓝图——被拉伸用于
飞行、划水、奔跑与书写：\omterm{def:g9:heredity-and-traits:heredity}{遗传}是重新分配功能，
而非重新设计。”陈列柜：鲸的髋骨、封闭的
甲虫翅膀与人类的尾骨。

\textbf{3.} 由字母计数隐含的树——
黑猩猩最近、\omterm{ex:g6:producing-our-food:bread}{酵母}最远——必须与隔壁的
解剖学家与古生物学家的树相吻合：\omterm{def:g1:plants-around-us:tree}{分支}相同，
却是独立推导得出。

\textbf{4.} 蛾类：\omterm{def:g8:reproduction-and-environment:population}{种群}的体色在数十年内追随其
\omterm{def:g6:exploring-our-environment:components}{环境}而变化。雀类：喙的
分布随\omterm{def:g2:from-seed-to-plant:seed}{种子}供应而移动，经卡尺测量
——\omterm{def:g9:heredity-and-traits:heredity}{遗传}驱动的变化，正被实时观察。

\textbf{5.} 因为世系会\omterm{def:g1:plants-around-us:tree}{分支}：我们与今日的
猴子是同一分叉点上不同\omterm{def:g1:plants-around-us:tree}{分支}的表亲——
没有任何活着的\omterm{def:g6:exploring-our-environment:components}{生物}是谁的祖先，一位表亲的存在
并不否认\omterm{prop:g6:classification-kinship:kinship}{亲缘关系}，正如你的表亲
也不会否认你的\omterm{prop:g6:classification-kinship:kinship}{亲缘关系}。

\textbf{6.} \emph{Archaeopteryx}，介于
\omterm{prop:g4:classifying-animals:classes}{爬行动物}骨架与\omterm{prop:g3:sorting-living-things:groups}{鸟类}组件之间；有腿的\omterm{def:g9:evidence-of-evolution:fossils}{化石}鲸则
介于陆生\omterm{prop:g3:sorting-living-things:groups}{哺乳动物}与游泳者之间。两者在被发现之前，
就已被预言。

\textbf{7.} 在科学中，“理论”指的是经受检验、
能够组织一个领域的解释，而非猜测。这些展厅
所确立的，是变化与同源的历史事实，
由多条相互独立的证据共同证明；真正意义上的理论，
则是下一展厅所阐释的机制。

\textbf{8.} 若真有时代错乱——\omterm{prop:g3:sorting-living-things:groups}{哺乳动物}的\omterm{def:g3:bones-and-muscles:skeleton}{骨骼}
位于三叶虫层之下，并经证实——便会打破
该论断所要求的顺序。两百年来没有任何一次发掘
产出过此类\omterm{def:g9:evidence-of-evolution:fossils}{化石}，这一论断每天都在经受最严峻的检验，
在一切仍在挖掘的地方持续通过。

\textbf{9.} “隔壁所见的套盒，
从来不是为了归档方便：特征之所以层层嵌套，
是因为世系会\omterm{def:g1:plants-around-us:tree}{分支}——分类本身就是家谱树，
在人们明白它所描绘的内容之前便已绘成。”

\textbf{10.} “本展区的每件展品，以及此前的每位参观者，
都以同样的四个字母书写其构造，
并由同一套密码读取：\omterm{def:g1:plants-around-us:tree}{分支}各不相同，
而主干的语言从未改变。”

\textbf{11.} “本展区表明：\omterm{def:g5:biodiversity:species}{物种}会发生变化，
它们共享同一世系，\omterm{def:g9:evidence-of-evolution:fossils}{化石}记录保存了
过渡类型。它尚未回答的，是什么\emph{驱动}
这种变化——下一展厅即为机制展厅，其答案
依赖于变异、\omterm{def:g9:heredity-and-traits:heredity}{遗传}与\omterm{def:g6:exploring-our-environment:components}{环境}的
制约。”

\textbf{12.} “一切生命是同一本家庭相册，
岩石保存了最早的几页——假以时日，表亲
会变得如鲸与蝙蝠般迥异，却仍是一家人。”
\end{solution}
